\subsection{Kernels, Ranges, and Map Properties}
Throughout this subsection, $\mathcal{U},\mathcal{V},\mathcal{W}$ are vector spaces over the same fixed field $\mathbb{F}$.

\begin{definition} \label{def:kernel-and-range-linear-map} \hlt{Kernel and Range of a Linear Map}\\
For $T\in\LL(\mathcal{V},\mathcal{W})$, define its \hlt{kernel}, or \hlt{nullspace}, by $\ker(T)\defeq\{\mathbf{v}\in\mathcal{V}:T\mathbf{v}=\mathbf{0}_{\mathcal{W}}\}$ and its range by $\operatorname{range}(T)\defeq\{T\mathbf{v}:\mathbf{v}\in\mathcal{V}\}$.
\end{definition}

\begin{proposition} \label{prop:injective-iff-kernel-zero} \hlt{Injectivity and the Kernel}\\
A linear map $T\in\LL(\mathcal{V},\mathcal{W})$ is injective iff $\ker(T)=\{\mathbf{0}_{\mathcal{V}}\}$.
\begin{proof}
\begin{enumerate}[label=\roman*.]
\setlength{\itemsep}{0pt}
\item[$(\Rightarrow)$] Suppose $T$ is injective. For $\mathbf{v}\in\ker(T)$, $T\mathbf{v}=\mathbf{0}_{\mathcal{W}}=T\mathbf{0}_{\mathcal{V}}$ by Proposition \ref{prop:linear-maps-send-zero-to-zero}, so $\mathbf{v}=\mathbf{0}_{\mathcal{V}}$ by injectivity.\\
Moreover, $\mathbf{0}_{\mathcal{V}}\in\ker(T)$ by the displayed equality; hence $\ker(T)=\{\mathbf{0}_{\mathcal{V}}\}$.
\item[$(\Leftarrow)$] Suppose $\ker(T)=\{\mathbf{0}_{\mathcal{V}}\}$. For $\mathbf{v},\mathbf{w}\in\mathcal{V}$ with $T\mathbf{v}=T\mathbf{w}$, $T(\mathbf{v}-\mathbf{w})=\mathbf{0}_{\mathcal{W}}$ by linearity, so $\mathbf{v}-\mathbf{w}\in\ker(T)$ and $\mathbf{v}=\mathbf{w}$.
\end{enumerate}
\hlt{Structural Commentary}: The forward direction compares kernel vectors with zero; the reverse moves equal images into the kernel.\\
\hlt{Key Insight}: Differences convert image collisions into kernel vectors.
\end{proof}
\end{proposition}

\begin{proposition} \label{prop:kernel-and-range-are-subspaces} \hlt{Kernel and Range Are Subspaces}\\
If $T\in\LL(\mathcal{V},\mathcal{W})$, then:
\begin{enumerate}[label=(\roman*)]
\setlength{\itemsep}{0pt}
\item $\ker(T)$ is a subspace of $\mathcal{V}$.
\item $\operatorname{range}(T)$ is a subspace of $\mathcal{W}$.
\end{enumerate}
\begin{proof}
\begin{enumerate}[label=(\roman*)]
\setlength{\itemsep}{0pt}
\item The kernel contains $\mathbf{0}_{\mathcal{V}}$ by Proposition \ref{prop:linear-maps-send-zero-to-zero}. For $\mathbf{u},\mathbf{v}\in\ker(T)$ and $a\in\mathbb{F}$, $T(a\mathbf{u}+\mathbf{v})=aT\mathbf{u}+T\mathbf{v}=\mathbf{0}_{\mathcal{W}}$, so $a\mathbf{u}+\mathbf{v}\in\ker(T)$.\\
Therefore $\ker(T)\leq\mathcal{V}$ by Theorem \ref{thm:subspacetest}.
\item The range contains $\mathbf{0}_{\mathcal{W}}=T\mathbf{0}_{\mathcal{V}}$ by Proposition \ref{prop:linear-maps-send-zero-to-zero}. For $T\mathbf{u},T\mathbf{v}\in\operatorname{range}(T)$ and $a\in\mathbb{F}$, $aT\mathbf{u}+T\mathbf{v}=T(a\mathbf{u}+\mathbf{v})\in\operatorname{range}(T)$.\\
Therefore $\operatorname{range}(T)\leq\mathcal{W}$ by Theorem \ref{thm:subspacetest}.
\end{enumerate}
\hlt{Structural Commentary}: Each proof verifies nonemptiness and linear-combination closure before invoking the subspace test.\\
\hlt{Key Insight}: Linearity stabilizes the zero fiber and the image.
\end{proof}
\end{proposition}

\begin{definition} \label{def:invertible-linear-map-and-isomorphism} \hlt{Invertible Linear Map and Isomorphism}\\
A linear map $T\in\LL(\mathcal{V},\mathcal{W})$ is \hlt{invertible} iff $(\exists\,S\in\LL(\mathcal{W},\mathcal{V}))\,[ST=I_{\mathcal{V}}\land TS=I_{\mathcal{W}}]$.\\
Such an $S$ is an \hlt{inverse} of $T$. An invertible linear map is an \hlt{isomorphism}; if an isomorphism $\mathcal{V}\to\mathcal{W}$ exists, write $\mathcal{V}\cong\mathcal{W}$.
\end{definition}

\begin{proposition} \label{prop:inverse-linear-map-is-unique} \hlt{Inverse Is Unique}\\
If $T\in\LL(\mathcal{V},\mathcal{W})$ is invertible, then $(\exists!\,S\in\LL(\mathcal{W},\mathcal{V}))\,[ST=I_{\mathcal{V}}\land TS=I_{\mathcal{W}}]$. Denote this map by $T^{-1}$.
\begin{proof}
An inverse exists by invertibility of $T$. Suppose $S_1,S_2\in\LL(\mathcal{W},\mathcal{V})$ satisfy $S_1T=I_{\mathcal{V}}=S_2T$ and $TS_1=I_{\mathcal{W}}=TS_2$.\\
Then $S_1=S_1I_{\mathcal{W}}=S_1(TS_2)=(S_1T)S_2=I_{\mathcal{V}}S_2=S_2$ by Proposition \ref{prop:linear-map-product-algebra}(ii)--(iii) and the candidate inverse identities.\\
\hlt{Structural Commentary}: Existence is definitional; associativity compares any two inverse candidates.\\
\hlt{Key Insight}: A two-sided inverse has no freedom.
\end{proof}
\end{proposition}

\begin{proposition} \label{prop:inverse-of-product-linear-maps} \hlt{Inverse of a Product Reverses Order}\\
Let $T\in\LL(\mathcal{U},\mathcal{V})$ and $S\in\LL(\mathcal{V},\mathcal{W})$ be invertible.
\begin{enumerate}[label=(\roman*)]
\setlength{\itemsep}{0pt}
\item $ST$ is invertible.
\item $(ST)^{-1}=T^{-1}S^{-1}$.
\end{enumerate}
\begin{proof}
\begin{enumerate}[label=(\roman*)]
\setlength{\itemsep}{0pt}
\item Set $R\defeq T^{-1}S^{-1}$. Then $R\in\LL(\mathcal{W},\mathcal{U})$ by Proposition \ref{prop:linear-map-product-algebra}(i).\\
The identity $R(ST)=T^{-1}(S^{-1}S)T=T^{-1}T=I_{\mathcal{U}}$ holds by Proposition \ref{prop:linear-map-product-algebra}(ii)--(iii) and the inverse identities.\\
Similarly, $(ST)R=S(TT^{-1})S^{-1}=SS^{-1}=I_{\mathcal{W}}$ by Proposition \ref{prop:linear-map-product-algebra}(ii)--(iii) and the inverse identities. Hence $ST$ is invertible.
\item Since $R$ is a two-sided inverse of $ST$, $(ST)^{-1}=T^{-1}S^{-1}$ by Proposition \ref{prop:inverse-linear-map-is-unique}.
\end{enumerate}
\hlt{Structural Commentary}: The reversed product is checked against both inverse identities, then identified by uniqueness.\\
\hlt{Key Insight}: Composition order reverses under inversion.
\end{proof}
\end{proposition}

\begin{proposition} \label{prop:linear-map-invertible-iff-bijective} \hlt{Invertibility and Bijectivity}\\
For $T\in\LL(\mathcal{V},\mathcal{W})$, the map $T$ is invertible iff it is both injective and surjective.
\begin{proof}
\begin{enumerate}[label=\roman*.]
\setlength{\itemsep}{0pt}
\item[$(\Rightarrow)$] Let $S=T^{-1}$. If $\mathbf{v},\mathbf{w}\in\mathcal{V}$ and $T\mathbf{v}=T\mathbf{w}$, then $\mathbf{v}=(ST)\mathbf{v}=S(T\mathbf{v})=S(T\mathbf{w})=(ST)\mathbf{w}=\mathbf{w}$, so $T$ is injective.\\
If $\mathbf{y}\in\mathcal{W}$, then $\mathbf{y}=(TS)\mathbf{y}=T(S\mathbf{y})$, so $T$ is surjective.
\item[$(\Leftarrow)$] Suppose $T$ is injective and surjective. For $\mathbf{y}\in\mathcal{W}$, a unique $\mathbf{v}\in\mathcal{V}$ satisfies $T\mathbf{v}=\mathbf{y}$ by surjectivity and injectivity; define $S\mathbf{y}\defeq\mathbf{v}$.\\
For $a,b\in\mathbb{F}$ and $\mathbf{y}_1,\mathbf{y}_2\in\mathcal{W}$, $S(a\mathbf{y}_1+b\mathbf{y}_2)=aS\mathbf{y}_1+bS\mathbf{y}_2$ by injectivity because both sides have $T$-image $a\mathbf{y}_1+b\mathbf{y}_2$.\\
Hence $S\in\LL(\mathcal{W},\mathcal{V})$. By construction, $(\forall\,\mathbf{y}\in\mathcal{W})\,(TS)\mathbf{y}=\mathbf{y}$, so $TS=I_{\mathcal{W}}$ by Assumption \ref{assump:function-extensionality}.\\
For $\mathbf{v}\in\mathcal{V}$, $T((ST)\mathbf{v})=T\mathbf{v}$, whence $(ST)\mathbf{v}=\mathbf{v}$ by injectivity. Therefore $ST=I_{\mathcal{V}}$ by Assumption \ref{assump:function-extensionality}, and $T$ is invertible.
\end{enumerate}
\hlt{Structural Commentary}: Surjectivity constructs the reverse map; injectivity makes it linear and enforces both inverse identities.\\
\hlt{Key Insight}: Unique preimages turn a bijection into a linear inverse.
\end{proof}
\end{proposition}

\begin{definition} \label{def:linear-map-morphism-terminology-vector-spaces} \hlt{Morphism Terminology for Linear Maps}\\
Within vector spaces, the following terminology applies to $T\in\LL(\mathcal{V},\mathcal{W})$:
\begin{enumerate}[label=(\roman*)]
\setlength{\itemsep}{0pt}
\item $T$ is a \hlt{homomorphism}.
\item $T$ is an \hlt{endomorphism} iff $\mathcal{V}=\mathcal{W}$.
\item $T$ is a \hlt{monomorphism}, or \hlt{embedding}, iff $T$ is injective; equivalently, $\ker(T)=\{\mathbf{0}_{\mathcal{V}}\}$ by Proposition \ref{prop:injective-iff-kernel-zero}.
\item $T$ is an \hlt{epimorphism} iff $T$ is surjective, equivalently $\operatorname{range}(T)=\mathcal{W}$.
\item $T$ is an \hlt{isomorphism} iff $T$ is invertible; equivalently, $T$ is injective and surjective by Proposition \ref{prop:linear-map-invertible-iff-bijective}.
\item $T$ is an \hlt{automorphism} iff $\mathcal{V}=\mathcal{W}$ and $T$ is invertible.
\end{enumerate}
\end{definition}
